\section{Both triangle inequalities, with exact polynomial cores}
\label{sec:triangles}

Let \(a,b,c>0\) be side lengths satisfying the strict triangle
inequalities. Write \(\ell=a+b+c\), \(s=\ell/2\), and
\[
 x=s-a,\qquad y=s-b,\qquad z=s-c.
\]
Then \(x,y,z>0\), \(a=y+z\), \(b=z+x\), \(c=x+y\),
and \(s=x+y+z\). Put \(T\geq0\) for the area. Heron's identity is
\begin{equation}\label{eq:heron-new}
 T^2=sxyz.
\end{equation}
For completeness, place a side of length \(c\) on the horizontal axis.
The third vertex has horizontal coordinate
\(u=(b^2+c^2-a^2)/(2c)\) and height \(h\), where
\(h^2=b^2-u^2\). Since \(T=ch/2\), clearing denominators gives
\[
 16T^2=2(a^2b^2+b^2c^2+c^2a^2)-a^4-b^4-c^4
 =(a+b+c)(-a+b+c)(a-b+c)(a+b-c),
\]
which is \eqref{eq:heron-new}. Both equalities are polynomial identities.
Conversely these formulas construct the triangle using a positive square
root, so the finite arithmetic presentation is not vacuous.

\begin{theorem}[Triangle isoperimetry]
For every such triangle,
\[
 \ell^2\geq12\sqrt3\,T.
\]
Equality holds exactly for an equilateral triangle.
\end{theorem}
\begin{proof}
Finite arithmetic--geometric mean, proved in the graph section, gives
\(s^3=(x+y+z)^3\geq27xyz\). Multiplying by \(16s\) gives
\[
 \ell^4=16s^4\geq432sxyz=432T^2.
\]
Both sides of the desired unsquared inequality are nonnegative, so this
is equivalent to the assertion. Equality in arithmetic--geometric mean
forces \(x=y=z\), equivalently \(a=b=c\); conversely these values
give equality. For rational side inputs the squared assertion is a
single exact rational inequality.
\end{proof}

\begin{theorem}[The stronger product-of-sides inequality]
\label{thm:triangle-product}
For every nondegenerate triangle,
\[
 T\leq\frac{\sqrt3}{4}(abc)^{2/3}.
\]
Equality holds exactly for an equilateral triangle. A radical-free
equivalent assertion is
\begin{equation}\label{eq:triangle-product-polynomial}
 4096T^6\leq27(abc)^4.
\end{equation}
\end{theorem}
\begin{proof}
Set \(p=x+y+z\), \(q=xy+yz+zx\), \(r=xyz\), and
\(w=(x+y)(y+z)(z+x)=pq-r=abc\).
Arithmetic--geometric mean gives
\[
 q^3\geq27r^2,\qquad pq\geq9r.
\]
The second inequality can be checked without division: expand
\(pq=3r+(x^2y+y^2z+z^2x)+(xy^2+yz^2+zx^2)\).
Each parenthesized triple has product \(r^3\), hence has sum at
least \(3r\). These are the finite rational forms of the same mean
inequality; no logarithm or differentiation is involved.

It follows that \(9w=9pq-9r\geq8pq\), and
\(pq^4=(pq)q^3\geq243r^3\). Consequently
\[
 9^4w^4\geq8^4p^4q^4
 \geq4096\cdot243\,p^3r^3.
\]
Divide by \(243\), noting \(9^4/243=27\), and use
\(T^2=pr\). This proves \eqref{eq:triangle-product-polynomial}.
Taking nonnegative sixth roots gives the displayed geometric inequality.
Equality forces equality in \(q^3\geq27r^2\), so
\(xy=yz=zx\); positivity yields \(x=y=z\). The converse is direct.
\end{proof}

The stronger inequality indeed implies the perimeter inequality: finite
arithmetic--geometric mean gives \(abc\leq(\ell/3)^3\), and therefore
\[
 T\leq\frac{\sqrt3}{4}\left(\frac\ell3\right)^2
 =\frac{\ell^2}{12\sqrt3}.
\]
For a wholly polynomial implementation let
\(H=(a+b+c)(-a+b+c)(a-b+c)(a+b-c)=16T^2\).
The two targets become, respectively,
\[
 27H\leq\ell^4,\qquad H^3\leq27(abc)^4.
\]
These are the targets recorded in the Lean outline and the exact-arithmetic
verification script. Degenerate triangles have \(T=0\) and satisfy both
inequalities. The equality classifications above concern nondegenerate
triangles. In the stronger inequality, a degenerate triple of the form
\((t,t,0)\), up to permutation, also has equality. Degenerate equality
in the perimeter inequality occurs only at the all-zero triple.

